\documentclass{article}
\usepackage[T1]{fontenc}
\usepackage{lmodern}
\usepackage{graphicx}
\usepackage{amsmath,amssymb}
\usepackage[a4paper,margin=2cm]{geometry}
\usepackage[absolute,overlay]{textpos}
\setlength{\TPHorizModule}{1mm}
\setlength{\TPVertModule}{1mm}
\usepackage[indonesian]{babel}
\usepackage{hyperref}
\setlength{\emergencystretch}{2em}
% unit: Halaman 3

\begin{document}

% === Halaman 3 (python/native) ===

Hybrid Cryptography

• Alice dan Bob sepakat mengenkripsi dan mendekripsi pesan dengan algoritma 

kriptografi simetri (misalkan AES). 

• Bagaimana cara Bob dapat mengetahui kunci AES yang digunakan oleh Alice (kunci 

AES)?

• Solusinya adalah dengan menggunakan hybrid cryptography.

• Hybrid cryptography: menggabungkan kriptografi kunci-simetri (misalkan AES) 

dengan kriptografi kunci-publik (misalkan RSA).

• Caranya sebagai berikut:

    1) Alice memiliki sepasang kunci privat (SKAlice) dan kunci publik (PKAlice) miliknya. 

    2) Bob juga memiliki sepasang kunci privat (SKBob) dan kunci publik (PKBob) miliknya.

Rinaldi M/II4020 Kriptografi

3

\end{document}
